\documentclass[10pt,a4paper]{amsart}
\usepackage[margin=19mm]{geometry}
\usepackage{amsmath,amssymb,mathtools}
\usepackage[scr=rsfs,cal=euler]{mathalfa}
\usepackage{xfrac}
\usepackage[unicode,hidelinks,bookmarks=false]{hyperref}
\newcommand{\ie}{i.e.\ }
\newcommand{\cf}{cf.\ }
\newcommand{\resp}{respectively\ }
\newcommand{\Subsection}{\subsection}
\providecommand{\nobreakdash}{\nobreak\hbox{-}}
\setlength{\emergencystretch}{2em}
\multlinegap0pt
\usepackage{bm}
\usepackage{bbm}
\usepackage{dsfont}
\usepackage{xspace}
\usepackage{cancel}
\usepackage{mathtools}
\usepackage{amsthm}
\makeatletter
\@ifundefined{lem}{\newtheorem{lem}{Lemma}}{}
\@ifundefined{prop}{\newtheorem{prop}{Proposition}}{}
\makeatother
\title[Sub-wavelength resonances in two-dimensional multi-layer ela]{Sub-wavelength resonances in two-dimensional multi-layer elastic media}
\author{Yan Jiang, Hongyu Liu, Fanbo Sun,, Yajuan Wang}
\date{Source: arXiv:2601.12821v1 (CC BY 4.0)}
\begin{document}
\maketitle

\noindent\emph{Source and licence.} Sub-wavelength resonances in two-dimensional multi-layer elastic media. Version arXiv:2601.12821v1 (CC BY 4.0), distributed under the Creative Commons Attribution 4.0 International licence, as recorded in the arXiv licence metadata for this exact version and shown on the abstract page https://arxiv.org/abs/2601.12821v1. Full rights notice: licenses/arxiv-2601.12821-CC-BY-4.0.txt.\par\smallskip

\noindent\textit{Editorial scope.} Counted unit: the Lem whose source label is \texttt{thmdiskomega} in the article (source0.tex lines 1410--1429, source label \texttt{thmdiskomega}), delivered together with the article's own complete original proof of it (source0.tex lines 1430--1465, one \texttt{proof} environment of 36 lines). No mathematics is written, repaired or re-derived: statement and proof are the article's own text line for line. Required context retained uncounted: lamea (lines 343--345); lameb (lines 347--349); hats1 (lines 378--380); hats2 (lines 402--404); matrices p and m (lines 743--753); disk singlelayer (lines 1358--1367). Every definition, display and label of those lines is retained unchanged. Omitted from this package and not redistributed: 1--342, 346--346, 350--377, 381--401, 405--742, 754--1357, 1368--1409, 1466--1881. Nothing inside a retained excerpt is omitted or altered: every retained body is delivered exactly as the article's lines are written, so no omission receipt of any kind accompanies this package. Citations: the retained excerpts carry no citation command at all, so no bibliography entry of the article is transcribed and no reference is redistributed. Reference adaptation: none was needed and none was made; every reference inside the delivered unit resolves inside it, so no unresolved reference or citation is produced. Printed slips kept literal: no printed word, symbol, hypothesis or number of the counted statement or of its proof was altered. Layout and class: the article's own class is \texttt{amsart} with packages including booktabs, multirow, graphicx, amssymb, epstopdf, cite, geometry, amsmath, amsfonts, amsthm, mathrsfs, dsfont, color, bm; the wrapper is the lane's pinned \texttt{amsart} editorial wrapper at 10pt on A4, which restates the theorem-like environments the retained text needs and adds the article's own macro definitions that the retained text needs (none), with extra wrapper packages (bm, bbm, dsfont, xspace, cancel, mathtools). The delivered numbering is the unit's own. Assets: no retained span contains a figure environment, a graphics-inclusion command, a PDF-page inclusion, a table or a TikZ picture; no image file is copied into this package, the package declares no asset dependency and no image file is redistributed with it. Licence: version arXiv:2601.12821v1 of this article is distributed under the Creative Commons Attribution 4.0 International licence, as recorded in the arXiv licence metadata for this exact version and shown on the abstract page https://arxiv.org/abs/2601.12821v1. A full rights notice is delivered at licenses/arxiv-2601.12821-CC-BY-4.0.txt. Characters: every retained excerpt is pure ASCII, so the delivered engine, XeLaTeX with its default Latin Modern text font, reports no missing character.
\begin{equation}\label{lamea}
a_{\lambda,\mu}:=(2\lambda+6\mu)\beta_2+(3\lambda+5\mu)\beta_3,
\end{equation}

\begin{equation}\label{lameb}
b_{\lambda,\mu}:=\beta_{1}(2\lambda+6\mu)+\beta_{2}(\lambda+5\mu)+\beta_{3}(\lambda+\mu)+\beta_{4}(3\lambda+5\mu),
\end{equation}

\begin{equation}\label{hats1}
\hat{\mathbf{S}}_{\partial D}^{\omega}[\boldsymbol{\varphi}](\mathbf{x}):=\mathbf{S}_{\partial D}[\boldsymbol{\varphi}](\mathbf{x})+\gamma_{\omega}\int_{\partial D}\boldsymbol{\varphi}(\mathbf{y})d s(\mathbf{y}),
\end{equation}

\begin{equation}\label{hats2}
\hat{\mathbf{S}}_{\partial D}^{\tau\omega}[\boldsymbol{\varphi}](\mathbf{x}):=\mathbf{S}_{\partial D}[\boldsymbol{\varphi}](\mathbf{x})+\gamma_{\tau\omega}\int_{\partial D}\boldsymbol{\varphi}(\mathbf{y})d s(\mathbf{y}).
\end{equation}

\begin{prop}\label{matrices p and m}
The elements of the matrices $\mathbf{P}$ and $\mathbf{M}$ are given by
$$
P_{ij}=a_{\lambda,\mu}\int_{ \partial D}\boldsymbol{\zeta}_i(\mathbf{x})ds(\mathbf{x})\cdot \int_{ D}\boldsymbol{\xi}_j(\mathbf{x})d\mathbf{x},\quad\\
$$
$$
M_{ij}={b}_{\lambda,\mu}\int_{ \partial D}\boldsymbol{\zeta}_i(\mathbf{x})ds(\mathbf{x}) \cdot\int_{ D}\boldsymbol{\xi}_j(\mathbf{x})d\mathbf{x}
-\int_{ D} {\mathbf{S}}_{\partial D}[\boldsymbol{\zeta}_i](\mathbf{x}) \cdot \boldsymbol{\xi}_j(\mathbf{x}) d\mathbf{x},
$$
where the constants ${a}_{\lambda,\mu}$ and ${b}_{\lambda,\mu}$ are given by \eqref{lamea} and \eqref{lameb}, respectively.
\end{prop}

\begin{lem}\label{disk singlelayer}
Let $D$ be a disk of radius $R$. For $\mathbf{x}\in D$, there holds
\begin{equation}\label{disk singlelayer1}
{\mathbf{S}}_{\partial D}[\boldsymbol{\zeta}_i](\mathbf{x})=\int_{\partial D} \boldsymbol{\Gamma}(\mathbf{x}-\mathbf{y}) \boldsymbol{\zeta}_i(\mathbf{y}) d s(\mathbf{y})=\left(\sigma_1R\ln R-\frac{\sigma_2R}{2}\right)\boldsymbol{\zeta}_i(\mathbf{x}),\ i=1,2.
\end{equation}
and
\begin{equation}\label{disk singlelayer2}
{\mathbf{S}}_{\partial D}[\boldsymbol{\zeta}_3](\mathbf{x})=\int_{\partial D} \boldsymbol{\Gamma}(\mathbf{x}-\mathbf{y}) \boldsymbol{\zeta}_3(\mathbf{y}) d s(\mathbf{y})={-\frac{R}{2\mu}}\boldsymbol{\zeta}_3(\mathbf{x}).
\end{equation}
\end{lem}

\begin{lem}\label{thmdiskomega}
Let $D$ be a disk of radius $R$. For $\delta\ll 1$, and $\epsilon\ll 1$, there exist three sub-wavelength resonance frequencies,
counted with their multiplicities, whose leading-order terms denoted can be determined by the equations as follows
\begin{equation}\label{formuladisk}
{\rho\omega^2\ln\omega p_{i}+\rho\omega^2\Big(\ln{(\sqrt{\rho}\tau)}p_{i}
+m_{i}\Big)-\epsilon q_{i}=0, \quad i=1,2,3.}%\nonumber
\end{equation}
where $p_{i}$, $m_{i}$ and $q_{i}(1 \leq i \leq 3)$ are given by
\begin{equation}\label{Peigenvalue}
p_{1}=p_{2}={a}_{\lambda,\mu}\pi R^2, \quad\quad  p_{3}=0,\nonumber
\end{equation}
\begin{equation}\label{Meigenvalue}
m_{1}=m_{2}={b}_{\lambda,\mu}\pi R^2-\frac{R^2}{2}\left(\sigma_{1}\ln{R}-\frac{\sigma_2}{2}\right),\quad\quad m_{3}=\frac{R^2}{8\mu},\nonumber
\end{equation}
and
\begin{equation}\label{Qeigenvalue}
q_{1}=q_{2}=\frac{\sigma_{1}\ln{R}-\frac{\sigma_2}{2}+2\pi \gamma_{\tau\omega}}{\sigma_{1}\ln{R}-\frac{\sigma_2}{2}+2\pi \gamma_{\omega}},  \quad\quad  q_{3}=1.\nonumber
\end{equation}
Here, the constants ${a}_{\lambda,\mu}$ and ${b}_{\lambda,\mu}$ are given by \eqref{lamea} and \eqref{lameb}, respectively.
\end{lem}

\begin{proof}
Consider first, that $\{\boldsymbol{\xi}_i\}_{i=1}^{3}$ and $\{\boldsymbol{\zeta}_i\}_{i=1}^{3}$ are orthonormal, it is clear that
\begin{equation}
P_{ij}=0,\quad \quad M_{ij}=0, \quad \quad Q_{ij}=0, \quad\quad\text{for} \ i\neq j. \nonumber
\end{equation}
Next, it follows from Proposition \ref{matrices p and m} and Lemma \ref{disk singlelayer} that $P_{33}=p_3$,
$$
P_{ii}=a_{\lambda,\mu}\int_{ \partial D}\boldsymbol{\zeta}_i(\mathbf{x})ds(\mathbf{x})\cdot \int_{ D}\boldsymbol{\xi}_i(\mathbf{x})d\mathbf{x}
=\frac{{a}_{\lambda,\mu}}{2\pi R}|\partial D|Vol(D)=p_i,\quad i=1,2,
$$
$$
M_{33}=-\int_{ D} {\mathbf{S}}_{\partial D}[\boldsymbol{\zeta}_3](\mathbf{x}) \cdot \boldsymbol{\xi}_3(\mathbf{x}) d\mathbf{x}
=\frac{R}{2\mu}\frac{1}{2\pi R^3}\int_{D}r^2=m_3,
$$
and
\begin{align}
M_{ii}&={b}_{\lambda,\mu}\int_{ \partial D}\boldsymbol{\zeta}_i(\mathbf{x})ds(\mathbf{x}) \cdot\int_{ D}\boldsymbol{\xi}_i(\mathbf{x})d\mathbf{x}
-\int_{ D} {\mathbf{S}}_{\partial D}[\boldsymbol{\zeta}_i](\mathbf{x}) \cdot \boldsymbol{\xi}_i(\mathbf{x}) d\mathbf{x}=m_i,\quad i=1,2.\nonumber
\end{align}
In addition, by using Lemma \ref{disk singlelayer} and equations \eqref{hats1}, \eqref{hats2}, we have
$$
\hat{\mathbf{S}}_{\partial D}^{{\tau\omega} }[\boldsymbol{\zeta}_i]=R\Big(\sigma_{1}\ln{R}-\frac{\sigma_2}{2}+2\pi \gamma_{\tau\omega}\Big)\boldsymbol{\zeta}_i,\quad i=1,2,
$$
$$
\hat{\mathbf{S}}_{\partial D}^{{\omega} }[\boldsymbol{\zeta}_i]=R\Big(\sigma_{1}\ln{R}-\frac{\sigma_2}{2}+2\pi \gamma_{\omega}\Big)\boldsymbol{\zeta}_i,\quad i=1,2,
$$
and
$$
\hat{\mathbf{S}}_{\partial D}^{{\tau\omega} }[\boldsymbol{\zeta}_3]=\hat{\mathbf{S}}_{\partial D}^{{\omega} }[\boldsymbol{\zeta}_3]={-\frac{R}{2\mu}}\boldsymbol{\zeta}_3.
$$
Hence,
\begin{align}
Q_{ii}&=\Big(\hat{\mathbf{S}}_{\partial D}^{{\tau\omega} }[\boldsymbol{\zeta}_i],(\hat{\mathbf{S}}_{\partial D}^{{\omega},* })^{-1}[\boldsymbol{\xi}_i]\Big)=R\Big(\sigma_{1}\ln{R}-\frac{\sigma_2}{2}+2\pi \gamma_{\omega}\Big)\Big((\hat{\mathbf{S}}_{\partial D}^{{\omega}})^{-1}[\boldsymbol{\zeta}_i],\boldsymbol{\xi}_i\Big)=q_i,\quad i=1,2.\nonumber
\end{align}
In the same way, we can obtain that $Q_{33}=q_3$. Finally, it reaches the proof.
\end{proof}

\end{document}
